
\section{格罗莫夫极限的唯一性}\label{sec:5-4}
\zhindex{geluomofujixiandeweiyixing}{格罗莫夫极限的唯一性}

本节专门证明下面的唯一性结论。

\begin{theorem}\label{thm:5-4-1}
设 $J^\nu\in\mathcal J_\tau(M,\omega)$ 是一列被 $\omega$ \BookRecordIndex{T0372}驯服的\BookRecordIndex{T0088}殆复结构，并在
$C^\infty$ 拓扑下收敛到 $J\in\mathcal J_\tau(M,\omega)$。设
$(u^\nu,z_1^\nu,\ldots,z_n^\nu)$ 是一列带 $n$ 个两两不同\BookRecordIndex{T0066}标记点的
$J^\nu$-全纯球面，并且\BookRecordIndex{T0141}格罗莫夫收敛到两个\BookRecordIndex{T0340}稳定映射
$(\mathbf u,\mathbf z)$ 和 $(\widetilde{\mathbf u},\widetilde{\mathbf z})$。则
$(\mathbf u,\mathbf z)$ 与 $(\widetilde{\mathbf u},\widetilde{\mathbf z})$ 等价。
\end{theorem}

为了证明这个结论，必须比较\BookRecordIndex{T0141}格罗莫夫收敛定义中出现的两组重标度映射
$\phi_\alpha^\nu$ 和 $\widetilde\phi_{\widetilde\alpha}^\nu$ 的行为。第一个引理说明：
若 $\alpha,\beta$ 是 $(\mathbf u,\mathbf z)$ 的树 $T$ 中两个相邻顶点，则任何重标度
$\psi^\nu$ 都不可能在它们之间产生一个非平凡\BookRecordIndex{T0275}气泡，因为没有足够的能量形成这样的
\BookRecordIndex{T0275}气泡。

\begin{lemma}\label{lem:5-4-2}
设 $J^\nu$ 如定理~\ref{thm:5-4-1} 所述，并设带\BookRecordIndex{T0066}标记点的 $J^\nu$-全纯球面列
$(u^\nu,z_1^\nu,\ldots,z_n^\nu)$ 通过重参数化序列 $\phi_\alpha^\nu\in G$
\BookRecordIndex{T0141}格罗莫夫收敛到\BookRecordIndex{T0340}稳定映射
$(\mathbf u,\mathbf z)\in\mathcal{SC}_{0,n}(M;J)$。再设 $\alpha E\beta$ 是一条边，
且 $\psi^\nu\in G$ 是满足以下条件的序列：
\begin{enumerate}[label=(\alph*),leftmargin=3em]
\item $(\phi_\alpha^\nu)^{-1}\circ\psi^\nu$ 在 $S^2\setminus\{w_0\}$ 上
\BookRecordIndex{T0416}逐紧一致收敛到 $z_{\alpha\beta}$；
\item $(\phi_\beta^\nu)^{-1}\circ\psi^\nu$ 在 $S^2\setminus\{w_1\}$ 上
\BookRecordIndex{T0416}逐紧一致收敛到 $z_{\beta\alpha}$。
\end{enumerate}
则 $u^\nu\circ\psi^\nu$ 在 $S^2\setminus\{w_0,w_1\}$ 上\BookRecordIndex{T0416}逐紧一致收敛到
$u_\alpha(z_{\alpha\beta})=u_\beta(z_{\beta\alpha})$，而且只要
$r<d_{S^2}(w_0,w_1)$，就有
\begin{equation}\label{eq:5-4-1}
  \begin{aligned}
  \lim_{\nu\to\infty}E\bigl(u^\nu\circ\psi^\nu;B_r(w_0)\bigr)
  &=m_{\beta\alpha}(\mathbf u),\\
  \lim_{\nu\to\infty}E\bigl(u^\nu\circ\psi^\nu;B_r(w_1)\bigr)
  &=m_{\alpha\beta}(\mathbf u).
  \end{aligned}
\end{equation}
此外，若 $\alpha_i\in T_{\alpha\beta}$，则
$(\psi^\nu)^{-1}(z_i^\nu)$ 收敛到 $w_1$；若
$\alpha_i\in T_{\beta\alpha}$，则它收敛到 $w_0$。
\end{lemma}

\begin{proof}
固定常数 $\varepsilon>0$。当 $\nu$ 充分大时，
\[
  \phi_{\beta\alpha}^\nu
  \bigl(S^2\setminus B_\varepsilon(z_{\alpha\beta})\bigr)
  \subset B_\varepsilon(z_{\beta\alpha}).
\]
两边作用 $\phi_\beta^\nu$，得到
\[
  \phi_\alpha^\nu\bigl(S^2\setminus B_\varepsilon(z_{\alpha\beta})\bigr)
  \cap
  \phi_\beta^\nu\bigl(S^2\setminus B_\varepsilon(z_{\beta\alpha})\bigr)
  =\varnothing.
\]
另一方面，由 (a)、(b)，当 $\nu$ 充分大时，
\[
  \begin{aligned}
  \phi_\alpha^\nu\bigl(S^2\setminus B_\varepsilon(z_{\alpha\beta})\bigr)
  &\subset\psi^\nu(B_\varepsilon(w_0)),\\
  \phi_\beta^\nu\bigl(S^2\setminus B_\varepsilon(z_{\beta\alpha})\bigr)
  &\subset\psi^\nu(B_\varepsilon(w_1)).
  \end{aligned}
\]
在这两个集合上积分形式 $(u^\nu)^*\omega$，得到
\[
  \begin{aligned}
  E\bigl(u^\nu\circ\psi^\nu;B_\varepsilon(w_0)\bigr)
  &\ge E\bigl(u_\alpha^\nu;S^2\setminus B_\varepsilon(z_{\alpha\beta})\bigr)
   =E(\mathbf u)-E\bigl(u_\alpha^\nu;B_\varepsilon(z_{\alpha\beta})\bigr),\\
  E\bigl(u^\nu\circ\psi^\nu;B_\varepsilon(w_1)\bigr)
  &\ge E\bigl(u_\beta^\nu;S^2\setminus B_\varepsilon(z_{\beta\alpha})\bigr)
   =E(\mathbf u)-E\bigl(u_\beta^\nu;B_\varepsilon(z_{\beta\alpha})\bigr).
  \end{aligned}
\]
对 $i=0,1$，记
\[
  m_\varepsilon(w_i)
  :=\liminf_{\nu\to\infty}
  E\bigl(u^\nu\circ\psi^\nu;B_\varepsilon(w_i)\bigr).
\]
函数 $\varepsilon\mapsto m_\varepsilon(w_i)$ 单调不减；由定义~\ref{def:5-2-1}
中的（能量）公理，
\[
  \begin{aligned}
  m(w_0):=\lim_{\varepsilon\to0}m_\varepsilon(w_0)
  &\ge E(\mathbf u)-m_{\alpha\beta}(\mathbf u)
   =m_{\beta\alpha}(\mathbf u),\\
  m(w_1):=\lim_{\varepsilon\to0}m_\varepsilon(w_1)
  &\ge E(\mathbf u)-m_{\beta\alpha}(\mathbf u)
   =m_{\alpha\beta}(\mathbf u).
  \end{aligned}
\]
若 $B_\varepsilon(w_0)\cap B_\varepsilon(w_1)=\varnothing$，则
\[
  m_{\beta\alpha}(\mathbf u)
  \le m_\varepsilon(w_0)
  \le E(\mathbf u)-m_\varepsilon(w_1)
  \le E(\mathbf u)-m_{\alpha\beta}(\mathbf u)
  =m_{\beta\alpha}(\mathbf u).
\]
这就\BookRecordIndex{T0326}推出式~\eqref{eq:5-4-1}，并说明 $u^\nu\circ\psi^\nu$ 在
$S^2\setminus\{w_0,w_1\}$ 上\BookRecordIndex{T0416}逐紧一致收敛到一个常值。下面证明这个常值等于
$u_\alpha(z_{\alpha\beta})=u_\beta(z_{\beta\alpha})$。任选
$z\in S^2\setminus\{w_0,w_1\}$，并考察序列
\[
  \zeta^\nu:=(\phi_\alpha^\nu)^{-1}\circ\psi^\nu(z).
\]
于是 $\zeta^\nu\to z_{\alpha\beta}$，且
$(\phi_{\alpha\beta}^\nu)^{-1}(\zeta^\nu)\to z_{\beta\alpha}$。因此，
命题~\ref{prop:4-7-2} 的假设成立，其中分别以
$u_\alpha^\nu,u_\beta^\nu,\phi_{\alpha\beta}^\nu,z_{\alpha\beta},
z_{\beta\alpha},u_\alpha,u_\beta$ 代替
$u^\nu,v^\nu,\psi^\nu,z_0,z_\infty,u,v$。所以序列
\[
  u^\nu\circ\psi^\nu(z)=u_\alpha^\nu(\zeta^\nu)
\]
收敛到 $u_\alpha(z_{\alpha\beta})$。

现在设 $\alpha_i\in T_{\alpha\beta}$。则 $z_{\alpha i}=z_{\alpha\beta}$，并且由
定理~\ref{thm:5-2-2}\textup{(iii)}，
\[
  z_{\beta i}
  =\lim_{\nu\to\infty}(\phi_\beta^\nu)^{-1}(z_i^\nu)
  \ne z_{\beta\alpha}.
\]
这说明 $(\psi^\nu)^{-1}(z_i^\nu)$ 收敛到 $w_1$；否则，取适当子列后，对这个序列
作用 \BookRecordIndex{T0040}Möbius 变换 $(\phi_\beta^\nu)^{-1}\circ\psi^\nu$，会得到极限
$z_{\beta\alpha}$。同理，只要 $\alpha_i\in T_{\beta\alpha}$，就有
$(\psi^\nu)^{-1}(z_i^\nu)\to w_0$。引理~\ref{lem:5-4-2} 得证。
\end{proof}

下一个引理说明：若 $\psi^\nu$ 是一个重标度，并且 $u^\nu\circ\psi^\nu$ 在某个
有限集之外收敛，则 $\psi^\nu$ 的尺度必定与某个 $\phi_\alpha^\nu$ 相当；这里的
含义是，复合序列 $(\phi_\alpha^\nu)^{-1}\circ\psi^\nu$ 有收敛子列。稳定性条件正是在
这里起关键作用。

\begin{lemma}\label{lem:5-4-3}
设 $J^\nu,J,(u^\nu,z_1^\nu,\ldots,z_n^\nu),(\mathbf u,\mathbf z)$ 和
$\phi_\alpha^\nu$ 如引理~\ref{lem:5-4-2} 所述。设 $\psi^\nu$ 是 $G$ 中的序列，
$Z\subset S^2$ 是有限集，且 $v:S^2\to M$ 是 $J$-全纯球面，并满足：
\begin{enumerate}[label=(\alph*),leftmargin=3em]
\item $v^\nu:=u^\nu\circ\psi^\nu$ 在 $S^2\setminus Z$ 上\BookRecordIndex{T0416}逐紧一致收敛到 $v$，
并在 $Z$ 处发生\BookRecordIndex{T0243}冒泡；
\item 对每个 $i$，极限
$\zeta_i=\lim_{\nu\to\infty}(\psi^\nu)^{-1}(z_i^\nu)$ 存在；
\item 若 $v$ 是常值映射，则 $\#Y\ge3$，其中
$Y:=Z\cup\{\zeta_i\mid1\le i\le n\}$。
\end{enumerate}
则存在顶点 $\alpha\in T$，使序列
$(\phi_\alpha^\nu)^{-1}\circ\psi^\nu$ 有收敛子列。
\end{lemma}

\begin{proof}
反设不存在这样的顶点，并任取 $\alpha_0\in T$。于是序列
\[
  \chi^\nu:=(\phi_{\alpha_0}^\nu)^{-1}\circ\psi^\nu\in G
\]
没有收敛子列。因此，由引理 D.1.2，存在 $w_0,z_0\in S^2$，使 $\chi^\nu$ 的一个
子列（仍记作 $\chi^\nu$）在 $S^2\setminus\{w_0\}$ 上\BookRecordIndex{T0416}逐紧一致收敛到 $z_0$。
我们断言
\begin{equation}\label{eq:5-4-2}
  z_0\in Z_{\alpha_0},
  \qquad w_0\in Z.
\end{equation}
反设 $z_0\notin Z_{\alpha_0}$。于是 $u_{\alpha_0}^\nu$ 在 $z_0$ 的一个邻域中一致
收敛到 $u_{\alpha_0}$。所以
$v^\nu=u_{\alpha_0}^\nu\circ\chi^\nu$ 在 $S^2\setminus\{w_0\}$ 上\BookRecordIndex{T0416}逐紧一致收敛到
常值 $u_{\alpha_0}(z_0)$。因而 $v$ 为常值，且 $Z\subset\{w_0\}$。此外，因为
$z_0\notin Z_{\alpha_0}$，至多有一个指标 $i$ 满足 $z_{\alpha_0i}=z_0$。若有这样的
$i$，则对 $j\ne i$，序列
$(\phi_{\alpha_0}^\nu)^{-1}(z_j^\nu)$ 收敛到
$z_{\alpha_0j}\ne z_0$，故
\[
  w_0
  =\lim_{\nu\to\infty}(\chi^\nu)^{-1}
     \bigl((\phi_{\alpha_0}^\nu)^{-1}(z_j^\nu)\bigr)
  =\lim_{\nu\to\infty}(\psi^\nu)^{-1}(z_j^\nu).
\]
因此 $Y\subset\{w_0,\zeta_i\}$。若不存在这样的 $i$，则
$Y=Z\subset\{w_0\}$。两种情形都与 (c) 矛盾。所以
$z_0\in Z_{\alpha_0}$；类似的论证表明 $w_0\in Z$。

由式~\eqref{eq:5-4-2}，存在一点 $w_0\in Z$ 和一个满足
$\alpha_0E\alpha_1$ 的顶点 $\alpha_1\in T$，使得必要时取子列后，
\[
  (\phi_{\alpha_0}^\nu)^{-1}\circ\psi^\nu
\]
在 $S^2\setminus\{w_0\}$ 上\BookRecordIndex{T0416}逐紧一致收敛到 $z_{\alpha_0\alpha_1}$。
序列 $(\phi_{\alpha_1}^\nu)^{-1}\circ\psi^\nu$ 仍没有收敛子列。于是用同样的论证，
再取子列后，可以假设存在
$z_{\alpha_1\alpha_2}\in Z_{\alpha_1}$ 和 $w_1\in Z$，使
\[
  (\phi_{\alpha_1}^\nu)^{-1}\circ\psi^\nu
\]
在 $S^2\setminus\{w_1\}$ 上\BookRecordIndex{T0416}逐紧一致收敛到 $z_{\alpha_1\alpha_2}$。
我们断言 $\alpha_0\ne\alpha_2$。否则，可以对满足 $\alpha_0E\alpha_1$ 的顶点
$\alpha=\alpha_0$、$\beta=\alpha_1$ 应用引理~\ref{lem:5-4-2}。那就会\BookRecordIndex{T0326}推出
$u^\nu\circ\psi^\nu$ 在 $S^2\setminus\{w_0,w_1\}$ 上\BookRecordIndex{T0416}逐紧一致收敛到常值
$u_{\alpha_0}(z_{\alpha_0\alpha_1})$，且对每个 $i$，
$(\psi^\nu)^{-1}(z_i^\nu)$ 收敛到 $w_0$ 或 $w_1$。但这会说明 $v$ 为常值且
$Z\subset\{w_0,w_1\}$，与 (c) 矛盾。所以 $\alpha_0\ne\alpha_2$。
归纳可得 $T$ 中的无限序列 $\alpha_0,\alpha_1,\alpha_2,\ldots$，满足对所有 $i$，
$\alpha_iE\alpha_{i+1}$ 且 $\alpha_i\ne\alpha_{i+2}$。在树中，这样的序列由两两不同的
顶点组成，与 $T$ 的有限性矛盾。引理~\ref{lem:5-4-3} 得证。
\end{proof}

\begin{proof}[定理~\ref{thm:5-4-1} 的证明]
设 $\{\phi_\alpha^\nu\}_{\alpha\in T}$ 和
$\{\widetilde\phi_{\widetilde\alpha}^\nu\}_{\widetilde\alpha\in\widetilde T}$
是\BookRecordIndex{T0141}格罗莫夫收敛定义中出现的两组序列。我们分五步证明
$(\mathbf u,\mathbf z)$ 与 $(\widetilde{\mathbf u},\widetilde{\mathbf z})$ 等价。
首先，把\BookRecordIndex{T0141}格罗莫夫收敛定义中的（重标度）条件推广到不由边连接的顶点对
$\alpha,\beta$。为此，把记号 $z_{\alpha\beta}$ 推广到任意一对不同顶点
$\alpha,\beta\in T$。具体地，若 $\alpha\ne\beta$，则有唯一一个与 $\alpha$ 相邻的
顶点 $\gamma\in T$，使得 $\beta$ 位于从 $\alpha$ 经 $\gamma$ 可达的子集
$T_{\alpha\gamma}$ 中。记
\[
  z_{\alpha\beta}:=z_{\alpha\gamma},
  \qquad \beta\in T_{\alpha\gamma}.
\]
对任意 $\alpha,\beta\in T$，还使用记号
\[
  \phi_{\alpha\beta}^\nu
  :=(\phi_\alpha^\nu)^{-1}\circ\phi_\beta^\nu.
\]

\medskip
\noindent\textbf{步骤 1.}\quad
\emph{若 $\alpha,\beta\in T$ 且 $\alpha\ne\beta$，则
$\phi_{\alpha\beta}^\nu$ 在 $S^2\setminus\{z_{\beta\alpha}\}$ 上\BookRecordIndex{T0416}逐紧一致收敛到
$z_{\alpha\beta}$。}

设 $\gamma_0,\ldots,\gamma_m$ 是 $T$ 中从 $\gamma_0=\alpha$ 到
$\gamma_m=\beta$ 的\BookRecordIndex{T0064}边链。由假设，$\phi_{\gamma_{i-1}\gamma_i}^\nu$ 在
$S^2\setminus\{z_{\gamma_i\gamma_{i-1}}\}$ 上\BookRecordIndex{T0416}逐紧一致收敛到
$z_{\gamma_{i-1}\gamma_i}$。由于 $\gamma_{i-2}\ne\gamma_i$，有
$z_{\gamma_{i-1}\gamma_i}\ne z_{\gamma_{i-1}\gamma_{i-2}}$。因此，
$\phi_{\gamma_{i-2}\gamma_i}^\nu$ 在
$S^2\setminus\{z_{\gamma_i\gamma_{i-2}}\}$ 上\BookRecordIndex{T0416}逐紧一致收敛到
$z_{\gamma_{i-2}\gamma_i}$。归纳可知，$\phi_{\gamma_0\gamma_i}^\nu$ 在
$S^2\setminus\{z_{\gamma_i\gamma_0}\}$ 上\BookRecordIndex{T0416}逐紧一致收敛到
$z_{\gamma_0\gamma_i}$。取 $i=m$，便得到步骤 1 的断言。

\medskip
\noindent\textbf{步骤 2.}\quad
\emph{存在唯一双射 $f:T\to\widetilde T$，以及一个仍以
$(u^\nu,z_i^\nu,\phi_\alpha^\nu)$ 表示的子列，使得对每个 $\alpha\in T$，极限}
\begin{equation}\label{eq:5-4-3}
  \psi_\alpha
  =\lim_{\nu\to\infty}
  (\phi_\alpha^\nu)^{-1}\circ\widetilde\phi_{f(\alpha)}^\nu
\end{equation}
\emph{都存在；这里的收敛在整个 $S^2$ 上一致。}

对序列 $\psi^\nu:=\widetilde\phi_{\widetilde\alpha}^\nu$ 应用
引理~\ref{lem:5-4-3}，得到顶点 $\alpha\in T$，使序列
$(\phi_\alpha^\nu)^{-1}\circ\widetilde\phi_{\widetilde\alpha}^\nu$ 有一致收敛的子列。
由步骤 1，$\alpha$ 是唯一的。现在，对这个子列以及任意另一个元素
$\widetilde\beta\in\widetilde T$ 再次应用引理~\ref{lem:5-4-3}，并归纳进行。
由此得到映射 $g:\widetilde T\to T$，以及一族 \BookRecordIndex{T0040}Möbius 变换
$\{\chi_{\widetilde\alpha}\}_{\widetilde\alpha\in\widetilde T}$，使得对每个
$\widetilde\alpha\in\widetilde T$，
\[
  \chi_{\widetilde\alpha}
  =\lim_{\nu\to\infty}
  (\phi_{g(\widetilde\alpha)}^\nu)^{-1}
  \circ\widetilde\phi_{\widetilde\alpha}^\nu.
\]
交换 $\mathbf u$ 与 $\widetilde{\mathbf u}$ 的角色，可知 $g$ 是双射。于是
$f:=g^{-1}:T\to\widetilde T$ 与 \BookRecordIndex{T0040}Möbius 变换
$\psi_\alpha:=\chi_{f(\alpha)}$ 满足步骤 2 的要求。

\medskip
\noindent\textbf{步骤 3.}\quad
\emph{设 $f$ 和 $\psi_\alpha$ 如步骤 2 所述。则}
\begin{equation}\label{eq:5-4-4}
  \widetilde u_{f(\alpha)}=u_\alpha\circ\psi_\alpha,
  \qquad
  \widetilde z_{f(\alpha)i}=\psi_\alpha^{-1}(z_{\alpha i}),
  \qquad
  \widetilde z_{f(\alpha)f(\beta)}=\psi_\alpha^{-1}(z_{\alpha\beta}),
\end{equation}
\emph{其中 $\alpha,\beta\in T$、$\alpha\ne\beta$，且 $i=1,\ldots,n$
\sourcefootnote{原书此处以及步骤 4 都把\BookRecordIndex{T0066}标记点指标的上界印作 $k$；但
定理~\ref{thm:5-4-1} 从头到尾固定的是 $n$ 个\BookRecordIndex{T0066}标记点，而本证明没有定义 $k$，故两处
统一改为 $n$。}。此外，若 $\alpha E\beta$，则序列
$u^\nu\circ\widetilde\phi_{f(\alpha)}^\nu$ 在点
$\widetilde z_{f(\alpha)f(\beta)}$ 处的\BookRecordIndex{T0276}气泡能量为}
\begin{equation}\label{eq:5-4-5}
  E\bigl(\{u^\nu\circ\widetilde\phi_{f(\alpha)}^\nu\};
    \widetilde z_{f(\alpha)f(\beta)}\bigr)
  =m_{\alpha\beta}(\mathbf u).
\end{equation}

式~\eqref{eq:5-4-4} 的第一个等式来自式~\eqref{eq:5-4-3}，即
\[
  u_\alpha\circ\psi_\alpha
  =\lim_{\nu\to\infty}u^\nu\circ\phi_\alpha^\nu\circ\psi_\alpha
  =\lim_{\nu\to\infty}u^\nu\circ\widetilde\phi_{f(\alpha)}^\nu
  =\widetilde u_{f(\alpha)}.
\]
这里的极限应理解为\BookRecordIndex{T0247}模冒泡收敛。式~\eqref{eq:5-4-4} 的第二个等式也用到
式~\eqref{eq:5-4-3}：
\[
  \psi_\alpha^{-1}(z_{\alpha i})
  =\psi_\alpha^{-1}\left(
    \lim_{\nu\to\infty}(\phi_\alpha^\nu)^{-1}(z_i^\nu)\right)
  =\lim_{\nu\to\infty}(\widetilde\phi_{f(\alpha)}^\nu)^{-1}(z_i^\nu)
  =\widetilde z_{f(\alpha)i}.
\]
为证明式~\eqref{eq:5-4-4} 的第三个等式，取不同的
$\alpha,\beta\in T$。由步骤 1，$\phi_{\alpha\beta}^\nu$ 在
$S^2\setminus\{z_{\beta\alpha}\}$ 上\BookRecordIndex{T0416}逐紧一致收敛到 $z_{\alpha\beta}$。因此，对
$z\ne z_{\beta\alpha}$，由式~\eqref{eq:5-4-3} 得
\[
  \begin{aligned}
  \psi_\alpha^{-1}(z_{\alpha\beta})
  &=\psi_\alpha^{-1}\left(
    \lim_{\nu\to\infty}\phi_{\alpha\beta}^\nu(z)\right)\\
  &=\lim_{\nu\to\infty}
    (\widetilde\phi_{f(\alpha)}^\nu)^{-1}\circ\phi_\beta^\nu(z)\\
  &=\lim_{\nu\to\infty}
    \widetilde\phi_{f(\alpha)f(\beta)}^\nu
    \bigl(\psi_\beta^{-1}(z)\bigr).
  \end{aligned}
\]
所以 $\widetilde\phi_{f(\alpha)f(\beta)}^\nu$ 在
$S^2\setminus\{\psi_\beta^{-1}(z_{\beta\alpha})\}$ 上\BookRecordIndex{T0416}逐紧一致收敛到
$\psi_\alpha^{-1}(z_{\alpha\beta})$。这就证明了式~\eqref{eq:5-4-4}。

为证明式~\eqref{eq:5-4-5}，注意函数
\[
  \varepsilon\longmapsto
  \lim_{\nu\to\infty}
  E\bigl(u^\nu\circ\widetilde\phi_{f(\alpha)}^\nu;
    B_\varepsilon(\widetilde z_{f(\alpha)f(\beta)})\bigr)
\]
在 $\varepsilon\to0$ 时的极限不会改变，如果把球
\[
  B_\varepsilon(\widetilde z_{f(\alpha)f(\beta)})
  =B_\varepsilon(\psi_\alpha^{-1}(z_{\alpha\beta}))
\]
换成集合 $\psi_\alpha^{-1}(B_\varepsilon(z_{\alpha\beta}))$。现在
\[
  \begin{aligned}
  &\lim_{\nu\to\infty}
  E\bigl(u^\nu\circ\widetilde\phi_{f(\alpha)}^\nu;
    \psi_\alpha^{-1}(B_\varepsilon(z_{\alpha\beta}))\bigr)\\
  &\quad=\lim_{\nu\to\infty}
  E\bigl(u^\nu\circ\widetilde\phi_{f(\alpha)}^\nu\circ\psi_\alpha^{-1};
    B_\varepsilon(z_{\alpha\beta})\bigr)\\
  &\quad=\lim_{\nu\to\infty}
  E\Bigl((u^\nu\circ\phi_\alpha^\nu)\circ
    \bigl((\phi_\alpha^\nu)^{-1}\circ
      \widetilde\phi_{f(\alpha)}^\nu\circ\psi_\alpha^{-1}\bigr);
    B_\varepsilon(z_{\alpha\beta})\Bigr)\\
  &\quad=\lim_{\nu\to\infty}
  E\bigl(u^\nu\circ\phi_\alpha^\nu;B_\varepsilon(z_{\alpha\beta})\bigr).
  \end{aligned}
\]
最后一个等号来自
\[
  \operatorname{id}
  =\lim_{\nu\to\infty}
  (\phi_\alpha^\nu)^{-1}\circ
  \widetilde\phi_{f(\alpha)}^\nu\circ\psi_\alpha^{-1}.
\]
再令 $\varepsilon\to0$，便得到式~\eqref{eq:5-4-5}。步骤 3 得证。

\medskip
\noindent\textbf{步骤 4.}\quad
\emph{设 $\alpha,\beta\in T$ 且 $\alpha\ne\beta$
\sourcefootnote{原书在此未写 $\alpha\ne\beta$；但上文只对不同顶点定义了
$z_{\alpha\beta}$，而边关系也只连接不同顶点，故补上这一条件。}。则
$\alpha E\beta$ 当且仅当}
\begin{equation}\label{eq:5-4-6}
  E\bigl(\{u_\alpha^\nu\};z_{\alpha\beta}\bigr)
  +E\bigl(\{u_\beta^\nu\};z_{\beta\alpha}\bigr)
  =E(\mathbf u)
\end{equation}
\emph{并且不存在 $i\in\{1,\ldots,n\}$ 同时满足
$z_{\alpha i}=z_{\alpha\beta}$ 和 $z_{\beta i}=z_{\beta\alpha}$。}

若 $\alpha E\beta$，则这些条件由\BookRecordIndex{T0141}格罗莫夫收敛的定义直接得到。反过来，设
$\alpha,\beta\in T$ 满足上述两个条件，并且对每个 $i$，都有
$z_{\alpha i}\ne z_{\alpha\beta}$ 或 $z_{\beta i}\ne z_{\beta\alpha}$。选取 $T$ 中
从 $\gamma_0=\alpha$ 到 $\gamma_m=\beta$ 的\BookRecordIndex{T0064}边链
$\gamma_0,\ldots,\gamma_m$。那么式~\eqref{eq:5-4-6} 等价于
\[
  m_{\gamma_0\gamma_1}(\mathbf u)
  +m_{\gamma_m\gamma_{m-1}}(\mathbf u)=E(\mathbf u).
\]
这\BookRecordIndex{T0326}推出
\[
  m_{\gamma_0\gamma_1}(\mathbf u)
  =m_{\gamma_{m-1}\gamma_m}(\mathbf u).
\]
若 $m\ne1$，则对每个
$\gamma\in T_{\gamma_0\gamma_1}\setminus T_{\gamma_{m-1}\gamma_m}$ 都有
$E(u_\gamma)=0$，而这个集合非空。因此，其中必有一个球面分支承载某个\BookRecordIndex{T0066}标记点
$z_i$。对这个 $i$，有
$z_{\alpha i}=z_{\alpha\beta}$ 且 $z_{\beta i}=z_{\beta\alpha}$，与假设矛盾。
所以 $m=1$，从而 $\alpha E\beta$。步骤 4 得证。

\medskip
\noindent\textbf{步骤 5.}\quad
\emph{$f$ 是\BookRecordIndex{T0304}树同构。}

步骤 3 表明，若 $\alpha E\beta$，则
\[
  E\bigl(\{u^\nu\circ\widetilde\phi_{f(\alpha)}^\nu\};
    \widetilde z_{f(\alpha)f(\beta)}\bigr)
  +E\bigl(\{u^\nu\circ\widetilde\phi_{f(\beta)}^\nu\};
    \widetilde z_{f(\beta)f(\alpha)}\bigr)
  =E(\widetilde{\mathbf u}).
\]
此外，对每个 $i$，都有
$\widetilde z_{f(\alpha)i}\ne\widetilde z_{f(\alpha)f(\beta)}$ 或
$\widetilde z_{f(\beta)i}\ne\widetilde z_{f(\beta)f(\alpha)}$。所以由步骤 4，
$f(\alpha)\,\widetilde E\,f(\beta)$。把 $f$ 换成 $f^{-1}$，可知
$f(\alpha)\,\widetilde E\,f(\beta)$ 又\BookRecordIndex{T0326}推出 $\alpha E\beta$。所以 $f$ 是\BookRecordIndex{T0304}树同构。
由步骤 5 和步骤 3，$(\mathbf u,\mathbf z)$ 与
$(\widetilde{\mathbf u},\widetilde{\mathbf z})$ 等价。定理~\ref{thm:5-4-1} 得证。
\end{proof}
